% =====================================================================
\chapter{Introduction}
\thispagestyle{plain}
% =====================================================================

\section{Introduction}

Whenever a mobile phone places a call, plays a video or runs an app, its traffic reaches the Internet through a \emph{radio access network} (RAN). The RAN turns the radio signal at the antenna into packets, and packets back into radio signals. It also schedules thousands of users on shared spectrum every half millisecond. Over three generations the RAN has changed shape. It started as a single cabinet at the foot of each tower and has become a chain of specialised units joined by packet networks. How many cells an operator can deploy, and at what cost, now depends on the rate, latency and reliability of the links in this chain, and most of all on the first one, the \emph{fronthaul}. This project deals with managing the load on that link.

\subsection{The mobile transport network: fronthaul, midhaul and backhaul}

Figure~\ref{fig:transport} traces a user's data through a 5G network. On the mast sits the radio unit (RU). It contains the antennas and the radio-frequency chain, and in the O-RAN split it also takes the lower part of the physical layer. The distributed unit (DU) handles the time-critical baseband, which means the rest of the physical layer, the MAC scheduler and RLC. Higher protocols with looser timing (PDCP, RRC, SDAP) run in the central unit (CU), and one CU can serve many DUs. Mobility, sessions and the connection to the Internet are left to the 5G core.

The requirements on the three transport segments between these units differ a great deal.
\begin{itemize}
  \item \textbf{Fronthaul} (RU $\leftrightarrow$ DU) is the hardest of the three. It carries digitised radio samples (IQ samples) at tens of Gbit/s per cell, with a one-way latency budget of about 100~$\mu$s. A packet that arrives late is of no use to the radio.
  \item \textbf{Midhaul} (DU $\leftrightarrow$ CU) sees packet data that has already passed the MAC/RLC layers. Its rate is close to the user data rate and it can tolerate a few milliseconds of latency.
  \item \textbf{Backhaul} (CU $\leftrightarrow$ core) brings the combined user traffic of many cells to the core network, where latencies of about 10~ms or more are acceptable.
\end{itemize}

\begin{figure}[H]
\centering
\begin{tikzpicture}[font=\small,>=Stealth,
  unit/.style={draw,rounded corners,align=center,minimum height=1.25cm,minimum width=1.75cm,inner sep=3pt},
  lbl/.style={font=\footnotesize\bfseries,align=center}]
  % UE
  \node[unit,fill=gray!12,minimum width=1.2cm] (ue) at (0,0) {UE\\[-1pt]\scriptsize phones};
  % tower
  \begin{scope}[shift={(2.25,-0.62)}]
    \draw[thick] (0,0) -- (0.32,1.55) -- (0.64,0);
    \draw (0.08,0.4) -- (0.56,0.4); \draw (0.16,0.8) -- (0.48,0.8); \draw (0.24,1.2) -- (0.40,1.2);
    \draw[fill=iitblue!35] (0.22,1.55) rectangle (0.42,1.95);
    \draw[iitblue] (0.52,1.85) arc (-30:30:0.25); \draw[iitblue] (0.62,1.85) arc (-30:30:0.38);
  \end{scope}
  \node[unit,fill=iitblue!15] (ru) at (4.0,0) {O-RU\\[-1pt]\scriptsize RF + low PHY};
  \node[unit,fill=good!15] (du) at (7.3,0) {O-DU\\[-1pt]\scriptsize high PHY,\\[-4pt]\scriptsize MAC, RLC};
  \node[unit,fill=orange!18] (cu) at (10.2,0) {O-CU\\[-1pt]\scriptsize PDCP, RRC};
  \node[unit,fill=purple!12] (core) at (13.0,0) {5G Core\\[-1pt]\scriptsize + Internet};
  \draw[<->,dashed,thick] (ue) -- (2.35,0);
  \node[font=\scriptsize] at (1.2,0.35) {air};
  \draw[<->,line width=2.6pt,accent] (ru) -- (du);
  \draw[<->,line width=1.6pt,orange!80!black] (du) -- (cu);
  \draw[<->,line width=1.0pt,purple!70!black] (cu) -- (core);
  \node[lbl,text=accent] at (5.65,1.2) {FRONTHAUL};
  \node[lbl,text=orange!80!black] at (8.75,1.2) {MIDHAUL};
  \node[lbl,text=purple!70!black] at (11.6,1.2) {BACKHAUL};
  \node[font=\scriptsize,align=center,text=accent] at (5.65,-1.45) {IQ samples\\7--100+ Gbit/s\\$\approx$100\,$\mu$s};
  \node[font=\scriptsize,align=center,text=orange!80!black] at (8.75,-1.45) {packets\\$\approx$ user rate\\$\approx$1--5\,ms};
  \node[font=\scriptsize,align=center,text=purple!70!black] at (11.6,-1.45) {aggregated\\traffic\\$\approx$10\,ms};
  \draw[decorate,decoration={brace,amplitude=6pt,mirror},thick] (3.1,-2.25) -- (14.0,-2.25)
     node[midway,below=8pt,font=\footnotesize] {Mobile transport network (``X-haul'')};
\end{tikzpicture}
\caption{End-to-end path of user data in a disaggregated 5G RAN, with the three transport segments and their typical rates and latency budgets. The fronthaul, highlighted, is the subject of this project.}
\label{fig:transport}
\end{figure}

\subsection{From distributed RAN to Open RAN}

A classical \emph{distributed RAN} (D-RAN) keeps the whole baseband unit (BBU) in a cabinet at each site, so only backhaul leaves the site. \emph{Centralised RAN} (C-RAN) pulled the BBUs of many sites into a central pool, where they could share hardware and coordinate cells. The cost of this was a fronthaul link carrying raw radio samples over the CPRI protocol. Its rate is constant and proportional to bandwidth $\times$ antennas, and it stays the same whether or not anyone is using the cell. Once radios have 64 antennas for massive MIMO and carriers are 100~MHz wide, CPRI would need hundreds of Gbit/s per cell, which is not practical. So 5G split the baseband in a new place (Figure~\ref{fig:evolution}). Part of the physical layer went back into the radio unit. The fronthaul now uses packet-based eCPRI on ordinary Ethernet and carries frequency-domain samples only for the resource blocks that are actually scheduled.

\begin{figure}[H]
\centering
\begin{tikzpicture}[font=\footnotesize,>=Stealth,
  b/.style={draw,rounded corners,align=center,minimum height=0.75cm,inner sep=2pt},
  row/.style={font=\footnotesize\bfseries,anchor=east,align=right}]
  % D-RAN
  \node[row] at (-0.3,0) {D-RAN};
  \node[b,fill=iitblue!15,minimum width=1.0cm] (r1) at (0.6,0) {RRH};
  \node[b,fill=good!15,minimum width=1.6cm] (b1) at (2.5,0) {BBU (at site)};
  \node[b,fill=purple!12,minimum width=1.4cm] (c1) at (9.5,0) {Core};
  \draw[<->,thick] (r1) -- (b1);
  \draw[<->,thick,purple!70!black] (b1) -- node[above]{backhaul} (c1);
  % C-RAN
  \node[row] at (-0.3,-1.6) {C-RAN};
  \node[b,fill=iitblue!15,minimum width=1.0cm] (r2) at (0.6,-1.6) {RRH};
  \node[b,fill=good!15,minimum width=1.9cm] (b2) at (5.0,-1.6) {BBU pool};
  \node[b,fill=purple!12,minimum width=1.4cm] (c2) at (9.5,-1.6) {Core};
  \draw[<->,line width=2.4pt,accent] (r2) -- node[above=2pt,text=accent]{CPRI: constant, huge} (b2);
  \draw[<->,thick,purple!70!black] (b2) -- node[above]{backhaul} (c2);
  % O-RAN
  \node[row] at (-0.3,-3.2) {O-RAN};
  \node[b,fill=iitblue!15,minimum width=1.0cm] (r3) at (0.6,-3.2) {O-RU};
  \node[b,fill=good!15,minimum width=1.2cm] (d3) at (4.0,-3.2) {O-DU};
  \node[b,fill=orange!18,minimum width=1.2cm] (u3) at (6.6,-3.2) {O-CU};
  \node[b,fill=purple!12,minimum width=1.4cm] (c3) at (9.5,-3.2) {Core};
  \draw[<->,line width=1.6pt,accent] (r3) -- node[above=2pt,text=accent]{eCPRI 7-2x} node[below=2pt,text=accent,font=\scriptsize]{load-dependent} (d3);
  \draw[<->,thick,orange!80!black] (d3) -- node[above]{midhaul} (u3);
  \draw[<->,thick,purple!70!black] (u3) -- node[above]{backhaul} (c3);
\end{tikzpicture}
\caption{Evolution of the RAN architecture. In C-RAN the CPRI fronthaul has a constant, very high rate; in O-RAN split 7-2x it carries only scheduled resource blocks, so the load follows the traffic and can be statistically multiplexed.}
\label{fig:evolution}
\end{figure}

\subsection{Functional splits and the O-RAN 7-2x split}

3GPP TR 38.801 lists eight points where the protocol stack can be cut between a central and a remote unit (Figure~\ref{fig:splits}). A high split such as Option 2 gives a lower transport rate that follows the traffic, but it leaves the central unit less room to coordinate. A low split such as Option 8 (CPRI) centralises more. The price is a transport rate that is huge and constant. For Open RAN, the O-RAN Alliance chose the intra-PHY split 7-2x as the fronthaul. The O-RU keeps the FFT, cyclic-prefix handling and, optionally, precoding. Channel estimation, equalisation, decoding and scheduling run in the O-DU. What crosses the fronthaul is frequency-domain IQ samples for every scheduled physical resource block (PRB) and every MIMO layer, compressed with \emph{block floating point} (BFP).

\begin{figure}[H]
\centering
\begin{tikzpicture}[font=\footnotesize,>=Stealth,
  layer/.style={draw,fill=soft,minimum width=2.1cm,minimum height=0.6cm,align=center},
  cut/.style={draw=accent,thick,dashed}]
  \node[layer] (rrc)  at (0,0)    {RRC};
  \node[layer] (pdcp) at (0,-0.6) {PDCP};
  \node[layer] (rlch) at (0,-1.2) {High RLC};
  \node[layer] (rlcl) at (0,-1.8) {Low RLC};
  \node[layer] (mach) at (0,-2.4) {High MAC};
  \node[layer] (macl) at (0,-3.0) {Low MAC};
  \node[layer] (phyh) at (0,-3.6) {High PHY};
  \node[layer] (phyl) at (0,-4.2) {Low PHY};
  \node[layer] (rf)   at (0,-4.8) {RF};
  \foreach \y/\o in {-0.3/1,-0.9/2,-1.5/3,-2.1/4,-2.7/5,-3.3/6,-3.9/7,-4.5/8}{
     \draw[cut] (-1.25,\y) -- (1.25,\y);
     \node[anchor=west,text=accent] at (1.3,\y) {Option \o};
  }
  \draw[line width=1.4pt,good] (-1.35,-3.9) -- (1.35,-3.9);
  \node[anchor=west,text=good,font=\footnotesize\bfseries] at (2.6,-3.9) {$\leftarrow$ O-RAN split 7-2x};
  \draw[->,thick] (-2.0,-0.1) -- (-2.0,-4.7);
  \node[rotate=90,font=\scriptsize,align=center] at (-2.5,-2.4) {lower split $\Rightarrow$ more centralisation,\\higher and less elastic fronthaul rate};
  \node[anchor=west,font=\scriptsize,align=left] at (6.1,-1.2) {\textbf{Option 2} (CU/DU, F1):\\packet rate $\approx$ user rate};
  \node[anchor=west,font=\scriptsize,align=left] at (6.1,-2.5) {\textbf{Option 7-2x} (DU/RU):\\frequency-domain IQ per scheduled PRB,\\BFP compressed, eCPRI over Ethernet};
  \node[anchor=west,font=\scriptsize,align=left] at (6.1,-4.4) {\textbf{Option 8} (CPRI):\\time-domain samples, constant rate};
\end{tikzpicture}
\caption{Functional split options of 3GPP TR 38.801 on the RAN protocol stack. The O-RAN fronthaul uses the intra-PHY split 7-2x.}
\label{fig:splits}
\end{figure}

\subsection{The O-RAN architecture}

O-RAN does more than open the fronthaul. It also adds two programmable controllers (Figure~\ref{fig:oran}). The non-real-time RIC lives inside the service management and orchestration (SMO) framework, where it trains models and runs ``rApps'' on time scales above one second. The near-real-time RIC runs ``xApps'', which control the O-CU and O-DU over the E2 interface in a 10~ms--1~s loop. A fronthaul budget controller like the one we build in Part~B fits here naturally. It can be deployed as an xApp or as a function inside the O-DU pool, and it works at the 10~ms time scale of the near-RT RIC.

\begin{figure}[H]
\centering
\begin{tikzpicture}[font=\footnotesize,>=Stealth,
  b/.style={draw,rounded corners,align=center,minimum height=0.85cm,inner sep=3pt}]
  \node[b,fill=purple!10,minimum width=10.5cm,minimum height=1.0cm] (smo) at (3.5,3.6) {Service Management and Orchestration (SMO)\quad $\supset$\quad \textbf{Non-RT RIC} (rApps, $>1$\,s)};
  \node[b,fill=orange!12,minimum width=7.6cm] (ric) at (3.5,2.1) {\textbf{Near-RT RIC} (xApps, 10\,ms -- 1\,s)\ \ e.g.\ \emph{fronthaul budget xApp}};
  \node[b,fill=orange!18,minimum width=2.0cm] (cu) at (0.6,0.4) {O-CU};
  \node[b,fill=good!15,minimum width=2.0cm] (du) at (3.5,0.4) {O-DU};
  \node[b,fill=iitblue!15,minimum width=2.0cm] (ru) at (7.3,0.4) {O-RU};
  \draw[<->,thick] (smo.south -| ric.north) -- node[right]{A1} (ric.north);
  \draw[<->,thick] (ric.south -| cu.north) -- node[left]{E2} (cu.north);
  \draw[<->,thick] (ric.south -| du.north) -- node[left]{E2} (du.north);
  \draw[<->,thick] (cu) -- node[above]{F1} (du);
  \draw[<->,line width=2.2pt,accent] (du) -- node[above,text=accent,font=\footnotesize\bfseries]{Open FH} node[below,text=accent,font=\scriptsize]{7-2x} (ru);
  \draw[<->,thick,dashed] (smo.east) -- ++(0.6,0) |- node[pos=0.75,above]{O1 / O2} (ru.east);
\end{tikzpicture}
\caption{O-RAN architecture with its open interfaces. The open fronthaul (7-2x) connects O-DU and O-RU; the near-RT RIC controls the RAN over E2 at the 10~ms time scale used by our Part~B controller.}
\label{fig:oran}
\end{figure}

\subsection{Why the fronthaul is loaded}

In a 7-2x cell, the fronthaul load depends on how many PRB-layers are scheduled. The number of user bits does not set it. For one PRB on one layer in one slot, the O-RAN BFP format sends 12 subcarriers $\times$ 14 symbols of $b$-bit I and Q mantissas, plus an exponent. With a 9-bit mantissa that comes to about 3{,}400 bits. At its peak, a 100~MHz carrier with 273 PRBs and 4 layers needs about 7.4~Gbit/s. A 64-antenna receiver that forwards every antenna stream needs about 15.4~Gbit/s, and DU-side precoding for 64 streams needs more than 118~Gbit/s. Cells rarely peak at the same moment. Operators use this by connecting many cells through one aggregation switch and one 25 or 100~Gbit Ethernet link, with a capacity well below the sum of the peaks (Figure~\ref{fig:shared}). Such \emph{statistical multiplexing} keeps Open RAN fronthaul affordable. But it also means the link fails when bursts happen together.

\begin{figure}[H]
\centering
\begin{tikzpicture}[font=\footnotesize,>=Stealth,
   ru/.style={draw,fill=iitblue!15,rounded corners,minimum width=1.7cm,minimum height=0.55cm,align=center},
   ll/.style={draw,fill=accent!15,rounded corners,minimum width=1.7cm,minimum height=0.55cm,align=center}]
  \node[ll] (r1) at (0,3.0) {O-RU 1 (LL)};
  \node[ll] (r2) at (0,2.3) {O-RU 2 (LL)};
  \node[ll] (r3) at (0,1.6) {O-RU 3 (LL)};
  \node[ru] (r4) at (0,0.9) {O-RU 4 (eMBB)};
  \node at (0,0.35) {$\vdots$};
  \node[ru] (r8) at (0,-0.3) {O-RU 8 (eMBB)};
  \node[draw,fill=gray!20,circle,minimum size=1.1cm,align=center] (sw) at (3.4,1.35) {switch};
  \node[draw,fill=good!15,rounded corners,minimum width=2.4cm,minimum height=1.4cm,align=center] (du) at (9.3,1.35) {O-DU pool\\\scriptsize + budget controller};
  \foreach \r in {r1,r2,r3,r4,r8} \draw[<->] (\r.east) -- (sw);
  \draw[<->,line width=3pt,accent] (sw) -- (du);
  \node[text=accent,align=center] at (6.1,1.95) {\textbf{one shared 25GE link}};
  \node[text=accent,align=center,font=\scriptsize] at (6.1,0.75) {$\sum$ cell peaks $= 2.44\times$ capacity};
  \node[align=left,font=\scriptsize,anchor=west] at (-1.0,-1.15) {LL: 2\,ms deadline, priority 3 \qquad eMBB: 10\,ms deadline, priority 1};
\end{tikzpicture}
\caption{Shared fronthaul scenario of Part~B: eight 7-2x cells with different deadlines share one over-subscribed Ethernet link through an aggregation switch.}
\label{fig:shared}
\end{figure}

\section{Motivation for BTP and Problem Statement}

Two bottlenecks limit how cheaply, and how densely, Open RAN cells can be deployed.
\begin{enumerate}
  \item \textbf{The data are large and the processing time is short.} O-RAN BFP compression does little more than shorten the numbers, and it gets about 2$\times$ over 16-bit I/Q. On top of that, the O-DU has to compute a precoder for all 273 PRBs in every 0.5~ms slot. Each PRB needs the inverse of a complex $K\times K$ Gram matrix. With a standard numpy + LAPACK implementation and 16 users this takes 3.2~ms, more than six times the slot length.
  \item \textbf{The link is shared and smaller than the peak.} If several cells burst at once, the shared link cannot carry all of it, and low-latency (LL) cells can wait only about 2~ms. Some coordinator has to decide beforehand how many bits each cell may send. It makes this decision every 10~ms using telemetry that is already 1~ms old, and it cannot see the next burst. Any cell that gets too small a share loses data.
\end{enumerate}

We tackle both bottlenecks with the same guiding principle: \emph{keep the mathematical structure of the problem, and learn only what is uncertain}.

\paragraph{Problem statement (Part A).} The first aim is IQ compression with error bounds that uses the low rank and delay-domain sparsity of the received signal to go well past the 2$\times$ of BFP. The second is a per-slot complex precoder that is fast and accurate enough to run inside the 0.5~ms slot on a general-purpose CPU.

\paragraph{Problem statement (Part B).} Suppose we have a forecast of each cell's fronthaul need over the next interval. We want a controller that sets per-cell budgets on a shared 7-2x link every 10~ms, before the traffic is seen, and that minimises the priority-weighted fraction of bits missing their deadline. We also want to show that it beats the strongest available controllers when all of them run on identical traffic and are compared with paired statistics.

\section{Objectives and Contributions}

\begin{figure}[H]
\centering
\begin{tikzpicture}[font=\footnotesize,>=Stealth,
   lever/.style={draw,rounded corners,align=left,text width=6.6cm,inner sep=6pt}]
  \node[lever,fill=iitblue!8] (A) at (0,0) {%
    \textbf{\textcolor{iitblue}{Part A: make the data smaller and the DU faster}}\\[3pt]
    $\bullet$ CSEE and RAS-BFP structure-aware IQ encoders\\
    $\bullet$ 8-method inversion benchmark, learned InverseNets\\
    $\bullet$ RMT-WS warm-start complex inversion\\
    $\bullet$ Fused C/AVX-512/OpenMP RZF precoder kernel\\
    $\bullet$ RU-side precoding: $-72\,\%$ downlink fronthaul};
  \node[lever,fill=accent!6] (B) at (7.4,0) {%
    \textbf{\textcolor{accent}{Part B: share one link well}}\\[3pt]
    $\bullet$ Deadline-feasible rate $\rstar$ as the demand target\\
    $\bullet$ B1: GBM quantiles + closed-form KKT rule\\
    $\bullet$ Theory: MDP view, Markov-fluid bounds on $\rstar$\\
    $\bullet$ B2: Bayesian link-aware Transformer (LCQT)\\
    $\bullet$ B3: back-pressure newsvendor (BPN)\\
    $\bullet$ 3 traffic families, 19+ controllers, paired tests};
  \node[draw,rounded corners,fill=good!12,text width=14cm,align=center,inner sep=5pt,anchor=north] (AB) at ($(B.south)+(-3.7,-0.7)$) {%
    Part A $\times$ Part B: the compression width sets \emph{how much} load reaches the link;
    the allocation rule decides \emph{who} loses data when bursts coincide.};
  \draw[->,thick] (A.south) -- (A.south |- AB.north);
  \draw[->,thick] (B.south) -- (B.south |- AB.north);
\end{tikzpicture}
\caption{The two levers of fronthaul load management developed in this project and how they combine.}
\label{fig:overview}
\end{figure}

This project makes the following contributions.
\begin{enumerate}
  \item \textbf{Structure-aware IQ compression.} We propose two encoders with error bounds for the reference-symbol IQ used in channel estimation. CSEE exploits delay-domain sparsity at $O(MN\log N)$ cost, and RAS-BFP uses a randomised SRHT sketch and QR at $O(rMN)$. In link-level checks, CSEE loses less than 0.05\,\% of the sum-rate at 31$\times$ and RAS-BFP less than 0.8\,\% at 4$\times$, where O-RAN BFP reaches about 2$\times$. At 64 antennas both encode faster than truncated SVD. Used on the reference symbols of each uplink slot, they save 4--20\,\% of the uplink fronthaul bits beyond what BFP gives.
  \item \textbf{Fast complex inversion for the precoder.} We screen 8 methods on CPU, including two learned Newton--Schulz inverters. We also develop RMT-WS, a Chebyshev/Newton--Schulz inversion guided by random-matrix theory. It tracks a warm start, follows a Doppler law and returns a residual certificate, and it needs 6 GEMMs where Neumann-type schemes need more than 10.
  \item \textbf{Kernel engineering.} Our fused C/AVX-512/OpenMP RZF precoder is 1.6--3.9$\times$ faster than a C precoder on OpenBLAS/LAPACK on one core, and 1.9--11$\times$ faster on four. Of the implementations we measured, it alone fits the whole precoder for 16 users into the 0.5~ms slot. That makes per-slot RU-side precoding possible, and with it the downlink fronthaul rate drops by 72\,\% while keeping 99.7\,\% of the exact sum-rate.
  \item \textbf{The KKT newsvendor rule (B1).} The rule forecasts the deadline-feasible rate $\rstar$ as a range and splits the link so that every cell carries equal priority-weighted risk. We prove four properties of it. In the NCC comparison it was the best of the 12 reference controllers in all scenarios of two traffic families (submitted to NCC 2027). Stronger baselines that we added later do beat it, and this is what led to B3.
  \item \textbf{Theory and stronger baselines.} We write the problem as a POMDP and derive Markov-fluid bounds linking $\rstar$ to effective bandwidth. We also add a third traffic family and the strongest baselines we could implement, namely predictive MaxWeight and cross-entropy policy search.
  \item \textbf{A Bayesian, link-aware Transformer forecaster (B2).} LCQT uses deadline-aware patch tokens with FiLM conditioning and cross-cell attention. Its outputs are non-crossing quantiles, and it ends in a Bayesian mixture layer followed by a differentiable KKT layer.
  \item \textbf{The back-pressure newsvendor (B3).} This is a two-stage policy class that includes both strongest baselines as exact special cases. Learned with the same search effort, it beats them in all 10 scenarios of three traffic families by 3--25\,\%. On held-out real 5G traces the margin is 6--7\,\%.
\end{enumerate}

\section{Organization of the Report}

Chapter~\ref{ch:lit} surveys prior work on fronthaul transport, IQ compression, massive-MIMO matrix inversion, predictive capacity allocation, uncertainty quantification and back-pressure control, and points out the gaps that this project addresses. The system model and methods of Part~A are given in Chapter~\ref{ch:parta-methods}, with their simulation results in Chapter~\ref{ch:parta-results}. Chapter~\ref{ch:partb-methods} covers Part~B. It sets up the system model and the deadline-feasible rate, and then develops the KKT newsvendor rule, the Bayesian Transformer forecaster and the back-pressure newsvendor. Their closed-loop simulation results on three traffic families appear in Chapter~\ref{ch:partb-results}. In Chapter~\ref{ch:impl} we describe the implementation and the software framework, along with the optimised CPU kernels and how they perform on real hardware. Chapter~\ref{ch:concl} concludes the report and outlines future work. Proofs, simulation parameters and instructions for reproducing the results are collected in the appendices.
